\documentclass[11pt]{article}

\usepackage{amsmath,amssymb}
\usepackage{booktabs}
\usepackage{graphicx}
\usepackage[margin=1in]{geometry}
\usepackage{hyperref}

\title{Two-Step Diffusion for KPP Equations: Sections 4.3 and 4.4}
\author{KPP Code Reproduction}
\date{\today}

\begin{document}
\maketitle

\section{KPP Front Speed Experiment on \(T^2\)}

This section follows the reference paper's Section~4.3. The goal is to
estimate the front speed of the KPP equation in a two-dimensional,
time-periodic cellular flow and to generate particle clouds that resemble the
long-horizon invariant measure of the underlying Feynman--Kac particle
system.

\subsection{Problem Setting and Feynman--Kac Particle System}

The reaction--diffusion--advection equation is

\begin{equation}
u_t+v(x)\cdot\nabla u=\kappa\Delta u+r u(1-u).
\end{equation}

Here \(\kappa\) is the diffusion constant, \(v\) is the prescribed velocity
field, and \(u\) is the reactant concentration. The minimal front speed along
a direction \(e\) is written as

\begin{equation}
c^*(e)=\inf_{\lambda>0}\frac{\mu(\lambda)}{\lambda},
\end{equation}

where \(\mu(\lambda)\) is the principal eigenvalue of a twisted parabolic
operator. For a fixed \(\lambda\), the corresponding twisted diffusion on
\(T^2=[0,2\pi)^2\) is

\begin{equation}
dX_t=\left(2\kappa\lambda e_1+v(X_t,t)\right)dt+\sqrt{2\kappa}\,dW_t.
\end{equation}

The cellular velocity field is

\begin{equation}
v(x,t)=\left(-\cos x_2-\theta\cos(2\pi t)\sin x_1,\,
\cos x_1+\theta\cos(2\pi t)\sin x_2\right).
\end{equation}

The Feynman--Kac weights use the potential

\begin{equation}
V_\lambda(x,t)=\kappa\lambda^2+\lambda v(x,t)\cdot e_1.
\end{equation}

This particle representation provides the reference invariant measures used
for training and evaluation.

\subsection{Experimental Setup}

The 2D experiment uses the parameters listed in Table~\ref{tab:2d_setup}.

\begin{table}[htbp]
\centering
\caption{2D KPP experimental parameters.}
\label{tab:2d_setup}
\begin{tabular}{lc}
\toprule
Parameter & Value \\
\midrule
Domain & \(T^2=[0,2\pi)^2\) \\
\(\lambda\) & \(2.0\) \\
\(\theta\) & \(1.0\) \\
Time step \(dt\) & \(2^{-8}\) \\
Substeps per period \(r_n\) & \(256\) \\
Particles per level \(M\) & \(2\times 10^4\) \\
Long-horizon generations \(T\) & \(4096\) \\
Cold start & uniform on \(T^2\) \\
Training \(\log_2\sigma\) & \(-2.00,\ldots,-3.75\) \\
Evaluation \(\log_2\sigma\) & \(-2.00,\ldots,-4.50\) \\
\bottomrule
\end{tabular}
\end{table}

Three initialization modes are compared: cold start, MeanFlow warm start, and
MeanFlow plus DP warm start. Two evaluation metrics are used. The first is
eigenvalue convergence of \(\hat c_T(\lambda)\) for
\(T=1,2,4,\ldots\). The second is the invariant-measure \(W_2\) on the torus.
In practice, the invariant-measure error is evaluated with a \(72\times 72\)
histogram and the periodic ground metric.

\subsection{Preliminary Results}

Table~\ref{tab:2d} reports current \(W_2\) values. The two-stage pipeline is
end-to-end functional. The DP stage reduces the error in several diffusion
regimes, but the current MeanFlow baseline is still higher than the reference
paper baseline.

\begin{table}[htbp]
\centering
\caption{2D KPP preliminary \(W_2\) results.}
\label{tab:2d}
\begin{tabular}{cccc}
\toprule
\(-\log_2\sigma\) & cold & MeanFlow & MF+DP \\
\midrule
2.00 & 1.054 & 0.256 & 0.123 \\
2.50 & 1.090 & 0.234 & 0.112 \\
3.00 & 1.114 & 0.221 & 0.125 \\
4.00 & 1.150 & 0.238 & 0.195 \\
4.50 & 1.172 & 0.236 & 0.261 \\
\bottomrule
\end{tabular}
\end{table}

\paragraph{Update interface.}
Replace these rows with the final \texttt{metrics.json} values after the
current 2D training completes.

\section{3D KPP Front Speed Experiment with Time-Dependent Kolmogorov Flow}

This section follows the reference paper's Section~4.4. The experiment extends
the two-stage approach to a three-dimensional, time-dependent Kolmogorov flow.

\subsection{Problem Setting}

The governing equation is the same KPP equation, but the spatial domain is
\(T^3=[0,2\pi)^3\). The velocity field is

\begin{equation}
v(x,t)=\left(
\sin(x_3+\sin(2\pi t)),
\sin(x_1+\sin(2\pi t)),
\sin(x_2+\sin(2\pi t))
\right).
\end{equation}

This flow is divergence-free and generates more complex anisotropic invariant
measures than the 2D cellular flow.

\subsection{Experimental Setup}

Table~\ref{tab:3d_setup} summarizes the 3D parameters. The discretization is
the same as in the 2D experiment.

\begin{table}[htbp]
\centering
\caption{3D KPP experimental parameters.}
\label{tab:3d_setup}
\begin{tabular}{lc}
\toprule
Parameter & Value \\
\midrule
Domain & \(T^3=[0,2\pi)^3\) \\
\(\lambda\) & \(2.0\) \\
Time step \(dt\) & \(2^{-8}\) \\
Substeps per period \(r_n\) & \(256\) \\
Particles per level \(M\) & \(2\times 10^4\) \\
Long-horizon generations \(T\) & \(4096\) \\
Training \(\log_2\sigma\) & \(-2.00,\ldots,-3.75\) \\
Evaluation \(\log_2\sigma\) & \(-2.00,\ldots,-4.50\) \\
\bottomrule
\end{tabular}
\end{table}

The invariant-measure error is computed on the \(x\)-\(z\) projection:

\begin{equation}
E_{xz}(\sigma)=W_2\left((\Pi_{xz})_\#\hat\mu(\sigma),
(\Pi_{xz})_\#\hat\nu(\sigma)\right),
\end{equation}

where \(\Pi_{xz}(x,y,z)=(x,z)\). Eigenvalue convergence is also reported for
warm and cold initializations.

\subsection{Preliminary Results}

Table~\ref{tab:3d} reports representative current \(W_2\) values. The 3D
MeanFlow baseline is closer to the paper than the 2D baseline, but the DP
improvement remains modest.

\begin{table}[htbp]
\centering
\caption{3D KPP preliminary \(W_2\) results.}
\label{tab:3d}
\begin{tabular}{cccc}
\toprule
\(-\log_2\sigma\) & cold & MeanFlow & MF+DP \\
\midrule
2.00 & 0.889 & 0.151 & 0.150 \\
2.50 & 0.933 & 0.136 & 0.135 \\
3.00 & 0.922 & 0.131 & 0.129 \\
4.00 & 0.957 & 0.133 & 0.132 \\
4.50 & 0.956 & 0.153 & 0.153 \\
\bottomrule
\end{tabular}
\end{table}

\paragraph{Update interface.}
After full 3D training, replace these values and record the actual seed from
\texttt{actual\_seed.json}.

\end{document}
